\section{Compact support and onto slow clocks}\label{ba:euclidean}

A whole-space processor can use arbitrarily separated finite lanes.
We shall construct a periodic cycle with fixed compact support and also
give a separate decaying force with the same material event. We first
establish the time change used for the latter choice. Its clock must tend
to infinity: traversing only a finite amount of machine time would not
preserve the halting event.

\subsection{The logarithmic clock at every derivative order}

\begin{lemma}[Logarithmic clock with full derivative bounds]
\label{ba:log-clock}
Let $W(s,x)$ be a smooth solenoidal field for $s\ge0$, initially zero,
with every mixed derivative bounded and support in one compact set $K$.
For fixed computable $\nu>0$, set
\[
 s(t)=\log(1+t),\quad \rho(t)=(1+t)^{-1},\quad
 U(t,x)=\rho(t)W(s(t),x),\quad
 F=U_t+(U\cdot\nabla)U-\nu\Delta U.
\]
Every derivative $\partial_t^k\partial_x^\alpha U$ and
$\partial_t^k\partial_x^\alpha F$ is uniformly
$O((1+t)^{-1-k})$ and belongs to space-time $L^2$.
Both supports lie in $K$.  If $\Xi$ is the material flow of $W$, that
of $U$ is $\Xi(s(t),a)$, so every fixed spatial reachability event is
unchanged.  Effective finite formulas for $W$ give effective formulas
for $U,F$ and all derivatives.
\end{lemma}
\begin{proof}\label{ba:log-clock-proof}
Since $s'=\rho$ and $\rho'=-\rho^2$,
\begin{equation}\label{ba:clock-rule}
 \partial_t\{\rho^jY(s(t),x)\}
       =\rho^{j+1}(\partial_s-j)Y(s(t),x).
\end{equation}
For $P_{k,j}(D)=\prod_{r=0}^{k-1}(D-j-r)$, with $P_{0,j}=1$,
induction yields
\begin{equation}\label{ba:clock-induction}
 \partial_t^k\{\rho^jY(s(t),x)\}
       =\rho^{j+k}P_{k,j}(\partial_s)Y(s(t),x).
\end{equation}
Spatial differentiation commutes with this identity.  Taking $j=1$
proves the velocity bound at every order.  The residual is exactly
\begin{equation}\label{ba:clock-residual}
 F=\rho^2\{W_s-W+(W\cdot\nabla)W\}(s(t),x)
       -\nu\rho\Delta W(s(t),x).
\end{equation}
Every mixed derivative of the braces is bounded by the finite product
rule and the hypotheses.  Apply \eqref{ba:clock-induction} with $j=2$
and $j=1$ to obtain orders $\rho^{2+k}$ and $\rho^{1+k}$,
respectively.  This gives the force bound, including every spatial
multi-index.  In particular no temporal derivative of the convection
term has been omitted.

Supports do not increase under differentiation or multiplication.
For either field $Z$, the squared integral of an indicated derivative
is at most
\[
 |K|C_{k,\alpha}^2\int_0^\infty(1+t)^{-2-2k}\,dt
       =\frac{|K|C_{k,\alpha}^2}{1+2k}.
\]
The zero datum remains zero, so residual realization applies.
The chain rule shows that $\Xi(s(t),a)$ solves the $U$ trajectory
equation with initial label $a$; uniqueness identifies it with the
actual trajectory.  The clock is increasing onto $[0,\infty)$ with
inverse $e^s-1$, so it loses no finite computational time.  Finally
logarithm on positive arguments, the rational powers of $1+t$, and
the finite-index evaluation of the template and its derivatives are
effective.  This proves all assertions.
\end{proof}


\subsection{Loaded Euclidean processors}

\begin{theorem}[Two full-box processors on Euclidean space]
\label{ba:euclidean-boxes}
For every finite machine and input, and every fixed computable $\nu>0$,
either the move-first recorder of Lemma~\ref{lem:bb-movefirst} or the
six-rule recorder of Lemma~\ref{ba:six-compiler} yields effective smooth forcing on $\mathbb R^3$ with one
compact spatial support, all mixed derivatives bounded, and period one
for $t\ge1$, such that
\[
 \exists t\ge0:\ (\Phi_t(0))_1<-1
 \quad\Longleftrightarrow\quad\text{halting}.
\]
The zero-data solution is unique in Lemma~\ref{ba:realization}'s class,
has pressure zero and uniformly bounded energy.  Separately, the six-rule
processor admits a force square-integrable on all space-time with the
same support and event.
\end{theorem}
\begin{proof}\label{ba:euclidean-boxes-proof}
Write $\Gamma$ for the full alphabet of the recorder being used.
For the six-rule table use initial frontier 2 and base $B=2|\Gamma|+2$.
Give nonhalting controls lanes $3i+[0,1]$ and halting controls lanes
$-3i+[0,1]$, enumerating each group from 1; the second coordinate is
$[0,1]$.  The minimal subdivision has gap $B^{-2}$.  Use full thickness
$[-1,1]$, heights $10i$, and the transported-cutoff construction in
Section~\ref{sec:ba-transported-localization}, with initial padding $1/(10B^3)$.
Target horizontal padding is at most $1/(10B^2)$, below half the gap;
vertical middle intervals are also separated.

For the move-first table use a single halt lane $-3+[0,1]$, other lanes
$2+2j+[0,1]$, base $B=2|\Gamma|+1$, the minimal subdivision, thickness
$[-1,1]$, and heights $4i$.  Use current-box padding $1/(10B^2)$.
The middle-phase separation is
$2(1+1/(10B^2))<4$; the planar gaps handle the other phases.
This localization has a fixed current padding, as distinct from the
transported initial padding in the six-rule version.

The rational initial code $P_*$ is reached from zero along
$\theta(t)P_*$ by a localized translation during $[0,1]$.
Repeating the chosen cycle produces a template $W$ with one compact
support and bounded mixed derivatives.  The code identity at times
$1+n$ follows by induction through the local rules; the respective
recorder invariant supplies every next checkpoint in finite positive
time.  A halt lies below $-1$ in the first coordinate.  In a nonhalting
run all branch endpoints have positive first coordinate, and
\eqref{eq:bb-convex-coordinate} keeps it positive throughout every cycle.
Loading is nonnegative.  Lemma~\ref{ba:realization} completes the PDE
argument.  Uniform support and boundedness give a uniform energy bound.
The separate square-integrable choice follows from
Lemma~\ref{ba:log-clock}, without claiming periodicity in physical
time for that choice.
\end{proof}


\subsection{Two realizations of a left frontier}

\begin{proposition}[Separate columns or solid-box parking]
\label{prop:bb-mid-stream}\label{prop:bb-mid-parking}
Both left-frontier tables have compact whole-space realizations,
periodic after $t=1$, with fixed label zero and event $X_1(t;0)<-1/2$.
The stopping version can use distinct-height stream functions and
has uniqueness with pointwise-in-time $H^1$ comparison pressure.
The continuing version transports full solid boxes by evacuation,
reshaping and delivery, and represents every local step, including
post-halt idle computation.  It has in particular uniqueness with
comparison pressure in $C([0,T];H^1)$.
\end{proposition}
\begin{proof}\label{proof:bb-mid-stream}\label{proof:bb-mid-parking}
Let $\Sigma$ be the full work--history--mark alphabet of the chosen
left-frontier table, including its record symbols and any halting idle
records in the continuing version. Use odd digits in base $B=2|\Sigma|+2$, offset $-2$ for $C_h$ and
$2,4,6,\ldots$ for other controls.  The two head-relative fractions
and height zero give a rational initial code $Z_*$.  The rightmost
nonempty history cell is the original cell $-2$, an absolute landmark.

For the stopping table, apply the stream-function construction
\eqref{eq:bb-mid-stream-path} at heights $z_i=3i$.
Choose source and target column paddings no greater than $1/2$ and
one third of their respective family gaps.  A middle layer padding
$1/2$ keeps distinct layers disjoint.  The explicit thickening bound
following that construction proves the neighborhood action.
Write $Z_*=(x_*,y_*,0)$. With a flat switch $\theta$ from zero to one,
load the origin along $\theta(t)Z_*$ using the velocity
$\nabla\times(0,0,H)$, where
$H(t,x,y,z)=\theta'(t)\chi(x,y,z)(x_*y-y_*x)$.
Choose $\chi$ compactly supported in $\R^3$ and equal to one near the
entire segment. The gate $\theta'$ vanishes on the time collars.
All nonhalt source and target lower-left first coordinates are at
least two.  Inequality~\eqref{eq:bb-mid-stream-safe} therefore
excludes the observer during each period; loading has nonnegative
first coordinate.  Halt codes lie in $(-2,-1)$.

For the continuing table retain its unrestricted first-row history
symbol and its idle work rules.  Thicken every rectangle by $[-1,1]$.
Choose $L$ ten larger than the maximum of zero and every source or
target right endpoint.  Reserve parking projections
$(L+4i,0)+[-1,1]^2$ and let $\delta$ be the minimum of one and the
gaps in each union of one rectangle family with these reservations.
Use padding $\delta/10$, travel height five, and $7N$ slots: evacuate
all sources by three translations each, reshape all parked boxes by
$\operatorname{diag}(1+\theta(\lambda_i-1),[1+\theta(\lambda_i-1)]^{-1},1)$,
and deliver by three translations each.  Horizontal widths stay at
most one, vertical half-width is one, and the reservations miss both
original families.  Lemma~\ref{lem:bb-lift-parking} proves exact action
on the whole boxes and their neighborhoods even with source--target
overlap.  A first localized translation loads $Z_*$ from the origin.
The sampling identity is now valid for every $k\ge0$, not just until
halting: $\Phi_{1+k}(0)=Z(C_k)$ for the continued table.

For a nonhalting run all used endpoint centers have first coordinate
at least two, parking centers are farther right, and half-widths are
at most $1/2$.  Translation, parking strain, and waiting therefore
keep first coordinate positive.  A halt checkpoint lies in $(-2,-1)$.
Proposition~\ref{prop:bb-realization} proves the analytic assertions
for both constructions.  Their compact support also gives uniformly
bounded kinetic energy.
\end{proof}
